\section{The Role of the Cabinet}

\subsection{Composition of the Cabinet}

The Regional Council brings together figures who occupy positions of political power across Latin America in November 1975. They represent different governments, institutions, and political traditions, ranging from military regimes to civilian administrations and influential figures within the armed forces and security apparatuses.

The members of the cabinet do not share a single ideology or political model. Their governments differ in their understanding of democracy, military authority, economic policy, and the role of the state. What they share is their position at the centre of their respective political systems and their interest in preserving the ability of those systems to determine their own futures.

This makes the cabinet inherently diverse. Developments in one country may be viewed as a warning, an opportunity, or a potential threat by another. Maintaining political relationships across these differences may therefore be as important as pursuing any particular ideological objective.

\subsection{Objectives}

The Regional Council is primarily concerned with maintaining the political position and stability of the governments it represents. This requires understanding an increasingly uncertain regional environment and assessing how developments within and beyond national borders could affect domestic authority.

This includes remaining informed about political opposition, revolutionary organizations, military developments, changes in neighboring governments, and the interests of external powers. Access to reliable information and the ability to respond to changing circumstances are therefore particularly valuable.

At the same time, political survival does not necessarily require permanent alignment with any particular actor. Maintaining diplomatic relations and remaining open to cooperation with different governments or external actors may aid when facing circumstances that could destabilize the grip on power.

\subsection{Internal tensions}

The diversity of interests of the individuals participating in the Council may overlap without being identical. Governments may have different assessments of the same events, while figures within the military and political establishments may hold differing views about the direction their countries should take.

These differences do not necessarily prevent diplomacy. However, the possibility of political change further complicates these relationships. The position occupied by a government or individual may not remain equally secure as circumstances develop. Consequently, maintaining a degree of flexibility may prove valuable, particularly when the intentions and capabilities of those around the table remain unclear.

\subsection{Relationship to other cabinets}

\subsubsection{Cabinet B - CIA}

The United States is an important external actor whose interests and capabilities may affect the political environment of Latin America. Its intelligence, diplomatic influence, and material resources may be valuable to some governments, while others may be more cautious about the implications of closer relations.

\subsubsection{Cabinet C - Revolutionaries}

Revolutionary organizations constitute a significant part of the political environment confronting the governments. Their objectives, capabilities, and relationships with one another vary considerably. How the cabinet distinguishes between armed revolutionary movements, democratic opposition, and other forms of dissent may become a major source of disagreement.

\subsubsection{Cabinet D - KGB}

The Soviet Union represents an external power whose interests may affect the balance between revolutionary movements and established governments. Its activity must therefore be assessed not only as an ideological threat, but according to how it could alter the political position and security of individual governments, and whether external aid is necessary to confront it.
